\section{共同基础设施}\label{sec:shared}

三套规则共享同一套数据与成本口径。这不是形式上的统一，而是论文结论本身的要求：论文所有比较的合法性都建立在“同一资产池、同一调仓频率、同一成本模型、同一评估指标、唯一区别是信号符号”这一设定上。若三套规则各自使用不同的成本假设，本文的规则之间就无法与论文的数字对照。

\subsection{数据口径}

\textbf{期货}。使用主力连续合约，并经换月跳空守门处理：先检测候选跳空（日收益超过前移滚动标准差的 10 倍且绝对值超过 2\%，且落在主力切换日 $\pm2$ 交易日窗口内），再用新主力合约当日收益或加权指数交叉确认，确认后做前复权调整。论文实证中的期货侧数据即按此处理。

\textbf{A 股}。通达信本地日线为不复权价，论文的修复规则是：主板个股单日收益低于 $-10.5\%$、创业板与科创板低于 $-20.5\%$ 者判定为除权，将该日收益置零（等价于按当日比例对更早价格做前复权）。论文样本内共识别并修复 $6121$ 处。

\textbf{必须补齐的缺口}。论文在审计章明确披露这项修复的缺陷：\emph{跳空幅度小于涨跌停幅度的除权无法识别}（如小额分红、$10$ 送 $1$），构成残余污染。本文的规则若用于实盘，必须补上这一层——可行做法是从数据商处取得复权因子序列，用它替换论文的近似规则。在补上之前，A 股反转规则的所有历史回测数字都带有这一已知偏差。

\textbf{指数}。A 股宽基指数仅用于两处：作为反转组合的对冲工具与业绩基准。论文已指出 7 个 A 股宽基指数的相关系数高达 $0.85$ 以上，其“组合”实质是单一 beta，不能用于分散。

\subsection{成本模型}

\begin{table}[htbp]
\centering\footnotesize
\caption{成本参数（与论文一致）}\label{tab:cost}
\begin{tabular}{@{}>{\raggedright\arraybackslash}p{4.6cm}>{\raggedright\arraybackslash}p{2.2cm}>{\raggedright\arraybackslash}p{7.5cm}@{}}
\toprule
\textbf{市场} & \textbf{单边成本} & \textbf{口径说明} \\
\midrule
期货 & 2 个基点 & 含手续费与滑点；论文口径 \\
A 股个股 & 15 个基点 & 含佣金、印花税（卖出千分之一）与滑点 \\
A 股指数/ETF & 5 个基点 & 对应 ETF 或股指期货实现 \\
\bottomrule
\end{tabular}
\end{table}

\textbf{成本计算规则}（沿用论文）：只对换入与换出部分收费，持有未动的仓位零成本。这在权重随收益漂移的设定下意味着：即使目标权重不变，实际权重已偏离，因此每次调仓仍会产生小额换手。

\textbf{大额资金必须升级成本模型}。论文的固定基点模型对反转类规则\emph{系统性低估}成本，原因是反转交易密集发生在流动性较差的时点（偏离越大时往往流动性越差，即流动性螺旋\cite{brunnermeier2009}）。实盘应改用参与率约束的冲击模型：

\begin{equation}
\text{冲击成本} = K\cdot\sqrt{\text{参与率}},\qquad
\text{参与率}=\frac{\text{本笔委托量}}{\text{同期市场成交量}} ,
\end{equation}

并设定参与率上限（论文工程约定为 $10\%$），超过上限的委托必须拆分或标记拒绝。$K$ 需按品种用历史成交数据标定。这一升级在论文中属定性讨论，本文将其列为\textbf{工程判断}。

\subsection{可交易性过滤}

三套规则共用同一套过滤，避免回测高估：

\begin{itemize}[leftmargin=2.2em]
  \item \textbf{涨跌停}：主板 $\pm10\%$，创业板（$300$/$301$）、科创板（$688$/$689$）$\pm20\%$。涨跌停当日无法成交，该标的的调仓顺延至下一可交易日。
  \item \textbf{停牌}：停牌期间不调仓；复牌首日不交易（价格发现未完成）。
  \item \textbf{流动性门槛}：期货要求近 20 日日均成交额不低于 1 亿元；A 股个股要求不低于 5000 万元。低于门槛者从候选池剔除。
  \item \textbf{次新与特殊处理}：A 股剔除 ST 与 *ST，剔除上市不足 250 个交易日的次新股。
\end{itemize}

论文披露其回测\emph{未做}涨跌停与停牌过滤，因此在高波动时段（如 2015 年）会高估可实现收益。本文把这些过滤写成规则的一部分，属于对论文局限的修补，标注为\textbf{工程判断}。

\subsection{评估指标}

三套规则统一报告 13 项指标，与论文一致：年化收益、年化波动、Sharpe、最大回撤、Calmar、偏度、超额峰度、年化换手、毛利年化、成本拖累年化、月胜率、Newey--West $t$ 统计量（5 阶滞后）、成本占毛利比。

其中三项在使用时需特别注意：

\textbf{成本占毛利比}是本文最看重的指标，因为它直接对应铁律二。论文中该比值从 $L=250$ 期货趋势的 $10.2\%$ 到 $L=5$ 期货趋势的 $50.1\%$ 不等；超过 $33\%$（即毛利不足成本 3 倍）的组合按铁律二废弃。

\textbf{偏度与超额峰度}用于判断策略的尾部特征。论文指出趋势策略总体右偏但含深度左尾（动量崩溃），反转策略左偏。规则上线前应检查实测偏度是否与预期方向一致；若趋势规则的实测偏度显著为负，说明它已退化为“没有凸性保护的追涨”。

\textbf{Newey--West $t$ 统计量}用于初步显著性判断。注意论文未做多重检验校正，因此 $t>2$ 只能视为“值得进一步检验”，不能视为“已证实显著”。论文中期货截面反转 $L=60$ 的 $t=2.59$ 是少数超过 2 的结果。
